%=============================================================================!
% 2.0  CHAPTER OPENING                                                        !
%=============================================================================!
Chapter~1 described motion through position, velocity and acceleration.
Once the acceleration is known at every instant, the starting position
and velocity fix the motion, but we have yet to establish what determines
that acceleration. Newton's laws connect it to the pushes and pulls on a
body, meaning the object whose motion we follow. We first treat the body
as a particle: a point that carries its mass, so that we can study its
motion from place to place before considering rotation in Chapter~5.

To use those pushes and pulls in a calculation, we need to measure them
as forces and establish the frames of reference in which the laws hold.
We then introduce mass and apply the laws to a falling ball, a bead sinking
through a thick liquid and a block on a spring. When several bodies
interact, adding their equations reveals a conserved quantity, the total
momentum, provided that the external forces sum to zero. The same
calculation gives the equations of molecular dynamics for a collection of
atoms. Chapter~3 develops a further consequence of these equations: the
conditions under which the total energy is conserved.
